\documentclass[10pt]{article}
\usepackage[T1]{fontenc}
\usepackage[utf8]{inputenc}
\usepackage[margin=0.72in]{geometry}
\usepackage{graphicx}
\usepackage{booktabs}
\usepackage{tabularx}
\usepackage{array}
\usepackage{xcolor}
\usepackage{hyperref}
\usepackage{float}
\usepackage[font=small,labelfont=bf]{caption}
\usepackage{enumitem}
\definecolor{navy}{HTML}{183B5B}
\definecolor{pale}{HTML}{EDF4F8}
\definecolor{warn}{HTML}{FFF4E5}
\hypersetup{colorlinks=true,linkcolor=navy,citecolor=navy,urlcolor=navy}
\setlength{\parindent}{0pt}
\setlength{\parskip}{0.45em}
\setlist[itemize]{nosep,leftmargin=1.4em}
\renewcommand{\arraystretch}{1.12}
\newcommand{\evidencebox}[2]{%
  \begin{center}\fcolorbox{navy}{pale}{\begin{minipage}{0.94\linewidth}
  \textbf{#1}\quad #2
  \end{minipage}}\end{center}}
\newcommand{\warningbox}[2]{%
  \begin{center}\fcolorbox{orange!75!black}{warn}{\begin{minipage}{0.94\linewidth}
  \textbf{#1}\quad #2
  \end{minipage}}\end{center}}
\newcommand{\metricdef}{Expert beats were paired one-to-one with detections inside $\pm75$ ms. TP is a paired detection, FP is an unpaired detector output, FN is an unpaired expert beat, and $F_1=2TP/(2TP+FP+FN)$.}
\newcommand{\provenance}{The ``historical'' rows reproduce the earlier NeuroKit/Pan--Tompkins hard-gate experiment. The ``current'' rows come from the complete all-48-record audit. Small NeuroKit count differences between audit versions reflect implementation/configuration differences and must not be interpreted as physiology.}

\title{MIT--BIH Record 113\\T-Wave Double-Detection Case Analysis}
\author{ECG-only seizure detector: RR--HRV branch audit}
\date{4 August 2026}

\begin{document}
\maketitle

\evidencebox{Bottom line.}{Record 113 is the clearest proof that detector failure is not the same as physical artifact. NeuroKit detected almost every expert beat but added 1,039 extra events on a record without a special noise warning. The old Pan--Tompkins veto solved this particular pattern, while the current UNSW primary detector achieved 1 FP and 1 FN. However, the present two-detector RR-support rule still marked most RR intervals unsupported because its NeuroKit secondary stream produced extra events.}

\section{Record and evidence status}
The MIT--BIH database records are approximately 30 minutes long, have two ECG channels sampled at 360 Hz, and include expert beat annotations \cite{moody2001mitbih,physionetmitdb,goldberger2000physiobank}. Record 113 is from a 24-year-old woman with MLII and V1 channels. The official summary describes predominantly normal sinus rhythm with rate variation possibly due to a wandering atrial pacemaker and gives no special noise warning \cite{physionetmitdb}.

\begin{tabularx}{\linewidth}{@{}>{\bfseries}p{0.22\linewidth}X@{}}
Observed & NeuroKit false candidates repeatedly lie on broad post-QRS waves in the displayed strip.\\
Measured & Their delay from the preceding expert QRS was approximately 306--356 ms in the audited example.\\
Inference & Timing and shape are consistent with T-wave double detection. The expert file does not label every T wave, so ``T wave'' is a supported interpretation, not annotation ground truth.\\
Rejected claim & These extra events do not prove motion artifact, electrode artifact, or seizure-related morphology.\\
\end{tabularx}

\section{Full-record results}
\metricdef\ \provenance

\begin{table}[H]
\centering
\caption{Record 113 full-record results.}
\begin{tabular}{@{}llrrrr@{}}
\toprule
Audit & Method & TP & FP & FN & $F_1$\\
\midrule
Historical & NeuroKit alone & 1,794 & 1,039 & 1 & 0.7753\\
Historical & NeuroKit + forward PT veto & 1,793 & 0 & 2 & 0.9994\\
Current & NeuroKit 300-ms baseline & 1,794 & 1,039 & 1 & 0.7753\\
Current & UNSW initial event & 1,794 & 1 & 1 & 0.9994\\
Current & UNSW refined R fiducial & 1,794 & 1 & 1 & 0.9994\\
\bottomrule
\end{tabular}
\end{table}

\begin{figure}[H]
\centering
\includegraphics[width=0.82\linewidth]{figures/record_113_current_metric_comparison.png}
\caption{Current full-record audit. NeuroKit's error is overwhelmingly false-positive double detection, not missed QRS complexes.}
\end{figure}

\section{Why the old Pan--Tompkins gate worked here}
Pan--Tompkins emphasizes rapid, high-energy QRS activity using band-pass filtering, a derivative, squaring, moving-window integration, and adaptive signal/noise thresholds \cite{pan1985qrs}. In this record its QRS event followed the genuine NeuroKit R-peak candidate within the tested 0--100-ms interval. The later broad wave occurred after that Pan--Tompkins QRS event had passed, so the false NeuroKit candidate received no new Pan--Tompkins support.

Consequently, 1,793 of 1,794 true NeuroKit detections passed while none of the 1,039 false detections passed. This is a highly successful \emph{record-specific} T-wave rejection result, not validation of the gate for other morphologies.

\begin{figure}[H]
\centering
\includegraphics[width=\linewidth]{figures/record_113_historical_hard_gate.png}
\caption{Historical experiment. Red crosses in the left panel recur on broad post-QRS waves; the Pan--Tompkins veto removes them in the right panel.}
\end{figure}

\clearpage
\section{The complete architecture exposes a new problem}
The Khamis/UNSW primary detector \cite{khamis2016telehealth,kristof2024qrs} removed the double-detection failure at the QRS-output level. Initial and refined outputs both had 1 FP and 1 FN. The displayed refined crosses align with expert QRS times.

\begin{figure}[H]
\centering
\includegraphics[width=\linewidth]{figures/record_113_complete_architecture_strip.png}
\caption{Current architecture. Purple ticks show the NeuroKit secondary stream. It often emits one event at the QRS and another on the following broad wave.}
\end{figure}

The two-detector RR rule was adapted from published work requiring support around both interval boundaries and the expected secondary-event count \cite{ho2024rrquality,ho2025rraccuracy}. Here only 634 of 1,794 RR intervals were supported (35.34\% coverage), despite primary-detector $F_1=0.9994$. All evaluable RR intervals had mean absolute error 1.37 ms, and the supported subset had 1.36 ms. Therefore the low coverage did \emph{not} identify a poor primary RR series; it mainly reflected contamination of the secondary detector by extra waves.

\warningbox{Critical architecture finding.}{A reliability rule is only as useful as its comparator. On record 113, a strong primary detector plus a weak secondary detector produces an over-conservative RR-quality label. The quality output must preserve its reason code (extra secondary events) rather than collapsing to ``artifact.''}

\section{Decision}
\begin{itemize}
\item Keep record 113 as the mandatory T-wave double-detection regression test.
\item Retain UNSW as the primary detector candidate.
\item Do not use NeuroKit secondary disagreement as a hard veto or as an artifact label.
\item Re-audit a comparator with explicit T-wave discrimination before using two-detector coverage to exclude HRV windows.
\end{itemize}

\begingroup\interlinepenalty=10000
\bibliographystyle{unsrt}
\bibliography{MITBIH_Record_Case_Reports}
\endgroup
\end{document}
